\documentclass[11pt, a4paper]{article}

\usepackage[utf8]{inputenc} % Required for inputting international characters
\usepackage[T1]{fontenc} % Output font encoding for international characters
% \usepackage{breqn}

\usepackage{nicefrac}
% \usepackage{units}
\usepackage{physics}

\usepackage{bbold}
\usepackage{graphicx}%
\usepackage{comment}
\usepackage{multicol}
\usepackage{booktabs}
\usepackage{xcolor}

% \usepackage{listings}
% \usepackage{upquote}
% \input{lst-color.sty}

\usepackage{minted}
\usepackage{microtype}
\usepackage{wasysym}
\usepackage{fancyvrb} % Monospace
\usepackage{xspace}   % xspace for new command definition
\usepackage{slashed}
\usepackage{environ}
\usepackage{framed}


\usepackage{pxfonts}

\usepackage{pdfpages}


\usepackage[backend=biber, 
           style=numeric-comp,
%            style=alphabetic,
           sorting=none]{biblatex}
\addbibresource{database.bib} % The filename of the bibliography

% \usepackage[autostyle=true]{csquotes} % Required to generate language-dependent quotes in the bibliography

\usepackage{pdfpages}

\usepackage{everyshi}
\usepackage{ifthen}

\newcommand{\ped}[1]{_{\scriptscriptstyle{\textrm{#1}}}}
\newcommand{\apx}[1]{^{\scriptscriptstyle{\textrm{#1}}}}
\newcommand{\peff}{p\ped{eff}}
\newcommand{\mlsp}{m\ped{LSP}}
\newcommand{\gstrong}{g\ped{s}}
\newcommand{\alphastrong}{\alpha\ped{s}}
\newcommand{\Weff}{W\ped{eff}}
\newcommand{\TeV}{{\rm TeV}}
% \newcommand{\MadDM}{\protect\Verb|MadDM|\xspace}
% \newcommand{\DarkSUSY}{\protect\Verb|DarkSUSY|\xspace}
% \newcommand{\MicrOMEGAs}{\protect\Verb|MicrOMEGAs|\xspace}
% \newcommand{\Madgraph}{\protect\Verb|Madgraph|\xspace}
% \newcommand{\CalcHEP}{\protect\Verb|CalcHEP|\xspace}
% \newcommand{\SuperIsoR}{\protect\Verb|SuperIso Relic|\xspace}
% \newcommand{\SuperIsoRF}{\protect\Verb|SuperIso Relic v4.1|\xspace}
% \newcommand{\SuperIso}{\protect\Verb|SuperIso|\xspace}
% \newcommand{\MARTY}{\protect\Verb|MARTY|\xspace}
% \newcommand{\FORTRAN}{\protect\Verb|FORTRAN|\xspace}
% \newcommand{\DarkPack}{\protect\Verb|DarkPACK|\xspace}
% \newcommand{\FormCalc}{\protect\Verb|FormCalc|\xspace}
\newcommand{\sigmav}{\expval{\sigma v}}

\newcommand{\critdens}{{\rho\apx{c}_0}}
\newcommand{\Ho}{H_0}
\newcommand{\zeq}{z\ped{eq}}
\newcommand{\Omk}{\Omega_{k}}
\newcommand{\Omr}{\Omega\ped{r}}
\newcommand{\Omm}{\Omega\ped{m}}
\newcommand{\Omb}{\Omega\ped{b}}
\newcommand{\Omc}{\Omega\ped{c}}
\newcommand{\Omfluid}{\Omega\ped{fluid}}
\newcommand{\Oml}{\Omega_{\Lambda}}
\newcommand{\Omg}{\Omega_{\gamma}}
\newcommand{\Omn}{\Omega_{\nu}}

\newcommand{\sgstar}{\sqrt{g\ped{$\ast$}}}
\newcommand{\heff}{h\ped{eff}}

\newcommand{\M}{\mathcal{M}}
\newcommand{\relphasespace}[1]{\frac{\dd[3] \vb{p}_{#1}}{(2 \pi)^3 2 E_{#1}}}
\newcommand{\1}{\mathbb{1}}
\renewcommand{\i}{\mathrm{i}}
\newcommand{\vmol}{{v_{\scriptscriptstyle\textrm{M\o l}}}}
\newcommand{\geff}{{g\ped{eff}}}
% \newcommand{\heff}{{h_{\rm eff}}}
\newcommand{\Tdi}{T_{d_i}}
\newcommand{\yeq}{{Y\ped{eq}}}
\newcommand{\sigmavmol}{\langle \sigma \vmol \rangle}
\newcommand{\gast}{{g_\ast^{1/2}}}
\newcommand{\sveff}{\sigmavmol\ped{eff}}

% \newcommand{\peff}{p\ped{eff}}
\newcommand{\GeV}{\mathrm{GeV}}

\newcommand{\Gf}{G\ped{F}}

\newcommand{\ee}{\mathrm{e}}
\newcommand{\ii}{\mathrm{i}}

\newcommand{\stilde}{\widetilde}
\newcommand{\sbar}[1]{{\bar #1}}
\newcommand{\Ld}{\mathcal{L}}
\newcommand{\cc}{\textrm{c.c.}}
\newcommand{\Hu}{H_u}
\newcommand{\Hd}{H_d}
\newcommand{\Hup}{{H_u^+}}
\newcommand{\Huo}{{H_u^0}}
\newcommand{\Hdo}{{H_d^0}}
\newcommand{\Hdm}{{H_d^-}}
\newcommand{\dg}[1]{{#1}^\dagger}

\newcommand{\yu}{{\vb{y_u}}}
\newcommand{\yd}{{\vb{y_d}}}
\newcommand{\ye}{{\vb{y_e}}}
\newcommand{\yl}{{\vb{y_l}}}

\newcommand{\gammaeq}{\Gamma\ped{eq}}

\DeclareMathAlphabet{\mathpzc}{OT1}{pzc}{m}{it}
% \renewcommand{\tilde}{\widetilde}
\newcommand{\U}{{\mathpzc{u}}}
\newcommand{\D}{{\mathpzc{d}}}
\newcommand{\E}{{\mathpzc{e}}}
% Definizione dell'ambiente cppblock in minted

% \usepackage{fvextra} % Required for changing the delimiter
% % Change the delimiter characters for code blocks
% \RecustomVerbatimEnvironment{Verbatim}{BVerbatim}{}

% \definecolor{bg}{rgb}{0.97,0.97,0.97}
% \definecolor{bg}{HTML}{FFFFF0}
\definecolor{bg}{HTML}{f0f1ea}
\definecolor{myred}{HTML}{ff0000} % Red
\definecolor{mygreen}{HTML}{00ff00} % Green
\definecolor{myblue}{HTML}{0000ff} % Blue
\definecolor{myblack}{HTML}{000000} % Black
\definecolor{myorange}{HTML}{ff7f00} % Orange

% Set the breaklines option for minted globally\setminted{breaklines=true}
% 
% \newminted[cppcode2]{cpp}{
%   style=customcpp
% }
% 
% 
% \newminted{codeblockcpp}{
%   commentcolor=mygreen,
%   keywordcolor=myblue,
%   stringcolor=myred,
%   numbercolor=myorange,
%   identifiercolor=myblack
% }

\newcommand{\verblb}[1]{\nolinkurl{#1}}%plain verbatim inline text with line breaks admitted

\usemintedstyle{borland}%bw vs borland xcode tango
% \usemintedstyle{customcpp}

\newenvironment{cppblock}{\VerbatimEnvironment\begin{minted}[breaklines, bgcolor=bg]{c++}}{\end{minted}}
\newenvironment{cppblocknumbered}[1]{\VerbatimEnvironment\begin{minted}[breaklines, linenos, firstnumber=#1, escapeinside=§§, bgcolor=bg]{c++}}{\end{minted}}
\newenvironment{cppblocknobg}{\VerbatimEnvironment\begin{minted}[breaklines]{c++}}{\end{minted}}

\newenvironment{mintedplain}{\VerbatimEnvironment\begin{minted}[breaklines]{text}}{\end{minted}}

\newcommand{\cppverb}[1]{\mintinline{c++}{#1}}
\newcommand{\plainverb}[1]{\mintinline{text}{#1}}

\newcommand{\MadDM}{\plainverb{MadDM}\xspace}
\newcommand{\DarkSUSY}{\plainverb{DarkSUSY}\xspace}
\newcommand{\MicrOMEGAs}{\plainverb{MicrOMEGAs}\xspace}
\newcommand{\Madgraph}{\plainverb{Madgraph}\xspace}
\newcommand{\CalcHEP}{\plainverb{CalcHEP}\xspace}
\newcommand{\SuperIsoR}{\plainverb{SuperIso Relic}\xspace}
\newcommand{\SuperIsoRF}{\plainverb{SuperIso Relic v4.1}\xspace}
\newcommand{\SuperIso}{\plainverb{SuperIso}\xspace}
\newcommand{\MARTY}{\plainverb{MARTY}\xspace}
\newcommand{\FORTRAN}{\plainverb{FORTRAN}\xspace}
\newcommand{\DarkPack}{\plainverb{DarkPACK}\xspace}
\newcommand{\FormCalc}{\plainverb{FormCalc}\xspace}
\newcommand{\Cpp}{\plainverb{C++}\xspace}


% Definizione di cppblock in listling
% \lstnewenvironment{cppblock}[1][]
% {\lstset{language=C++,style=Nblock,#1}}
% {}
% \newcommand{\cppverb}[1]{\lstinline{#1}}

\newcommand{\MSbar}{{\overline{\rm MS}}}

\newcommand{\rhoD}{\rho\ped{D}}
\newcommand{\rhorad}{\rho\ped{rad}}
\newcommand{\srad}{s\ped{rad}}
\newcommand{\sD}{s\ped{D}}
\newcommand{\dsrad}{\dot{s}\ped{rad}}
\newcommand{\dsD}{\dot{s}\ped{D}}
\newcommand{\Sigmarad}{\Sigma\ped{rad}}
\newcommand{\Sigmatstar}{\tilde\Sigma^\ast}
\newcommand{\SigmaD}{\Sigma\ped{D}}
\newcommand{\Mpl}{M\ped{Planck}}
\newcommand{\homosD}{s\ped{D,hom}}


\title{Thermodynamics in modified cosmological scenarios}
\begin{document}
\maketitle
\section{Introduction}
Let us consider a modified cosmological scenario, by the presence of a ``dark'' energy
density $\rhoD$, in the radiation epoch. The Hubble rate follows the equation:
\begin{align}
 H^2 = \frac{8 \pi G}{3}(\rhorad + \rhoD)
\end{align}
where $\rhorad$ reads, as usual:
\begin{align}
 \rhorad = \frac{\pi^2}{30}\geff(T) T^4
\end{align}
where $T$, unless specified otherwise, is the temperature of the radiation.
For the sake of later convenience, we express:
\begin{align}
 H = H\ped{$\Lambda$CDM}\tilde H
\end{align}
with $H\ped{$\Lambda$CDM}$ the Hubble parameter in the standard cosmological model:
\begin{align}
 H\ped{$\Lambda$CDM} = \sqrt{\frac{8 \pi G}{3}\rhorad} = \sqrt{\frac{4\pi^3}{45}\geff}\frac{1}{\Mpl}T^2
\end{align}
and $\tilde H$ parametrises the presence of dark energy:
\begin{align}
 \tilde H = \sqrt{1+ \frac{\rhoD}{\rhorad}}.
\end{align}
We also define:
\begin{align}
 \sqrt{g_{\ast}} = \frac{\heff}{\sqrt{\geff}} \qty(\frac{T}{3\heff}\dv{\heff}{T} +1),
\end{align}

Depending on the nature of the ``dark'' energy, there can be a modification of the entropy density, that can be written always as:
\begin{align}\label{eqn:totals}
 s = \srad + \sD
\end{align}
where the relation between $\srad$ and $T$ reads, as usual:
\begin{align}\label{expr:sradlcdm}
 \srad = \frac{2\pi^2}{45}\heff(T) T^3,
\end{align}
and whose derivative reads
\begin{align}\label{expr:dsraddT}
 \dv{\srad}{T} = 3 \frac{\srad}{T} \qty( \frac{T}{3\heff}\dv{\heff}{T} +1).
\end{align}


If $\sD =0$, then $s=\srad$, and $\srad$ follows the equation:
\begin{align}\label{eqn:sradlcdm}
 \dsrad = -3H \srad
\end{align}
where the dot denotes the total derivarive with respect to the cosmological time $t$.
By combining this last result with the expressions \eqref{expr:sradlcdm} and \eqref{expr:dsraddT}, we can derive the relation
between $t$ and $T$ in the standard cosmological model:
\begin{align}\label{eqn:dTdTlcdm}
 \dv{T}{t} = - \frac{H}{T}\qty( \frac{T}{3\heff}\dv{\heff}{T} +1).
\end{align}

\section{Modification of the entropy density}

We shall consider next 3 different modification to the entropy density, which may
overlap with each others.

\subsection{Standard entropy injection}
Let us now consider a modified cosmological scenario in which equation \eqref{eqn:sradlcdm}
is modified by adding an injection term:
\begin{align}\label{eqn:sradlcdminj}
 \dsrad = -3H \srad +\Sigmarad,
\end{align}
or, with a different parametrisation:
\begin{align}\label{eqn:sradlcdminj2}
 \dsrad = -3H\qty(1 - \Sigmatstar) \srad
\end{align}
In this case, we can keep $\srad(T)$ as in expression \eqref{expr:sradlcdm}, and,
consequently, its derivative with respect to $T$ as in expression \eqref{expr:dsraddT}; however,
the relation between $t$ and $T$ has to be re-determined,
since such a relation reparametrises correctly $\srad(t)$ with the additional
injection. By doing so, one obtains:
\begin{align}\label{eqn:dTdTsrad}
 \dv{T}{t} = - \frac{H\qty(1 - \Sigmatstar)}{T}\qty( \frac{T}{3\heff}\dv{\heff}{T} +1),
\end{align}
which correctly reduces to \eqref{eqn:dTdTlcdm} for $\Sigmatstar=0$.

In scenario we have just described, the modification of the entropy density involves
only the radiation, with an effect that affects the relation between $t$ ad $T$.
This scenario can be combined with the other two.

\subsection{Dark entropy production}
In this case, we suppose that we are in the case of equation \eqref{eqn:totals},
with $\sD$ known. We would like to parametrise $\dot s$ as:
\begin{align}\label{eqn:dots_dark}
 \dot s = -3Hs + \SigmaD + \Sigmarad,
\end{align}
and to determine $\SigmaD$ as a function of the temperature. Firstly, let us
write:
\begin{align}
 \SigmaD= \dot s + 3Hs -\Sigmarad
\end{align}
Now, let us expand $s = \srad + \sD$, and use equation \eqref{eqn:sradlcdminj2}:
\begin{align}
 \SigmaD &= \dsD + \dsrad + 3H(\srad + \sD) -\Sigmarad
\end{align}
Let us use equation \eqref{eqn:sradlcdminj} to eliminate $\dsrad$:
\begin{align}
 \SigmaD &= \dsD + 3H\sD =\\
         &= 3H\qty[\sD  + \frac{1}{3H}\dv{\sD}{T}\dv{T}{t}].
\end{align}
We can now use equation \eqref{eqn:dTdTsrad} to have:
\begin{align}\label{expr:SigmaD1}
 \SigmaD = 3H\qty[\sD - \frac{1 - \Sigmatstar}{\frac{T}{3\heff}\dv{\heff}{T} +1}\frac{T}{3}\dv{\sD}{T}].
\end{align}

This result can be further manipulated, to fully explicit the dependence on the
thermodinamical degrees of freedom, and remove the explicit dependence on $H$.
By using the definition of $\tilde H$ and $\sqrt{g_\ast}$, we can rewrite
\eqref{expr:SigmaD1} as:
\begin{align}\label{expr:SigmaD2}
 \SigmaD = \sqrt{\frac{4\pi^3}{5}}\frac{\tilde H}{\Mpl} T^2 \qty[\sqrt{\geff}\sD - \frac{T\qty(1 - \Sigmatstar)\heff}{3\sqrt{g_\ast}}\dv{\sD}{T}].
\end{align}
In terms of $\tilde H$, one can also write the relation among $\Sigmarad$
and $\Sigmatstar$ at a given temperature $T$:
\begin{align}
 \Sigmatstar = \frac{45\sqrt{5}\Mpl}{4 \pi^{7/2}\heff \sqrt{\geff} T^5 \tilde H}\Sigmarad.
\end{align}
Let us remark that equation \eqref{expr:SigmaD2} can be inverted to easily compute
$\dv{\sD}{T}$ given $\sD$ and $\SigmaD$.


\subsection{Dark entropy injection}
This case is analogous to the one of ``dark entropy production''. However,
instead of knowing the form of $\sD$, we know the form of $\SigmaD$, as it
appears in equation \eqref{eqn:dots_dark}. Then, we are interested in computing
$\sD$ from the known expression of $\SigmaD$.
In order to do that, we need to invert the relation \eqref{expr:SigmaD2}.
What we find is a first order linear differential equation:
The relation we find is the following:
\begin{align}
 \dv{\sD}{T} = \frac{3 \sqrt{\geff}\sqrt{g_\ast}}{T\qty(1 - \Sigmatstar)\heff }\sD - \sqrt{\frac{45}{4 \pi^3}}\Mpl \frac{\sqrt{g_\ast}}{\qty(1 - \Sigmatstar)\qty(\heff T^3) \tilde H}\SigmaD.
\end{align}
The solution to the associate homogeneous equation reads:
\begin{align}
 \homosD(T)= C_0 \exp(\int \dd T \frac{3 \sqrt{\geff}\sqrt{g_\ast}}{T\qty(1 - \Sigmatstar)\heff } )
\end{align}
We can work out the integrand, and find a more explicit form for $H$:
\begin{align}
 \frac{3 \sqrt{\geff}}{T\qty(1 - \Sigmatstar)\heff }&\frac{\heff}{\sqrt{\geff}} \qty(\frac{T}{3\heff}\dv{\heff}{T} +1)
  = \frac{3}{T\qty(1 - \Sigmatstar)}\qty(\frac{T}{3} \dv{\log(\heff)}{T} + 1)=\\
  &=\frac{1}{1 - \Sigmatstar}\qty(\dv{\log(\heff)}{T} + \dv{\log(T^3)}{T})=\\
  &=\frac{1}{1 - \Sigmatstar}\dv{\log(\heff T^3)}{T}
\end{align}
At this point, we can integrate by parts:
\begin{align}
 \int\dd T \frac{1}{1 - \Sigmatstar}\dv{\log(\heff T^3)}{T} = \frac{1}{1 - \Sigmatstar}\log(\heff T^3) -
 \int \dd T \log(\heff T^3)\frac{1}{(1-\Sigmatstar)^2}\dv{\Sigmatstar}{T}
\end{align}
and finally we obtain:
\begin{align}
 \homosD(T)= C_0 (\heff T^3)^\frac{1}{1 - \Sigmatstar}\exp(-\int \dd T \log(\heff T^3)\frac{1}{(1-\Sigmatstar)^2}\dv{\Sigmatstar}{T})
\end{align}
For the sake of convenience, we can define
\begin{align}
 \kappa(T) &= \qty(\heff(T) T^3)^\frac{1}{1 - \Sigmatstar(T)}\\
 \sigma_s(T) &= \exp(- \int_0^T \dd\theta \varepsilon(\theta)) \\
 \varepsilon(T) &= \dv{\Sigmatstar(T)}{T} \frac{\log(\heff(T) T^3)}{\qty(1 - \Sigmatstar(T))^2}
\end{align}
and have
\begin{align}
 \homosD(T)= C_0\kappa(T)\sigma_s(T)
\end{align}

Now we need to find the particular solution. We can put
the full equation in the form:
\begin{align}\label{eqn:diffsandJ}
 \dv{\sD}{T} + \frac{1}{J}\dv{J}{T}\sD=- \sqrt{\frac{45}{4 \pi^3}}\Mpl \frac{\sqrt{g_\ast}}{\qty(1 - \Sigmatstar)\qty(\heff T^3) \tilde H}\SigmaD,
\end{align}
where $J$ is a function of the temperature such that:
\begin{align}\label{eqn:diffJ}
 \frac{1}{J}\dv{J}{T}=-\frac{3 \sqrt{\geff}\sqrt{g_\ast}}{T\qty(1 - \Sigmatstar)\heff }.
\end{align}
In fact, once $J$ is known, we can write equation \eqref{eqn:diffsandJ} as:
\begin{align}
 \frac{1}{J}\dv{(J\sD)}{T} =- \sqrt{\frac{45}{4 \pi^3}}\Mpl \frac{\sqrt{g_\ast}}{\qty(1 - \Sigmatstar)\qty(\heff T^3) \tilde H}\SigmaD,
\end{align}
and therefore the form of a particular solution is:
\begin{equation}\label{eqn:particularsD}
   \sD(T) = - \frac{1}{J(T)}\qty{\sqrt{\frac{45}{4 \pi^3}}\Mpl \int^T_{0} \dd \theta
 \qty[\frac{\sqrt{g_\ast(\theta)}\SigmaD(\theta)J(\theta)}{\qty(1 - \Sigmatstar(\theta))\heff(\theta)\theta^3\tilde H(\theta)} ] - C_1}
\end{equation}
We are only left to find $J$. We can easily relate it to $\homosD$:
\begin{align}
 \frac{1}{J}\dv{J}{T}&=-\frac{3 \sqrt{\geff}\sqrt{g_\ast}}{T\qty(1 - \Sigmatstar)\heff } = -\frac{1}{\homosD}\dv{\homosD}{T} \\
  \dv{\log(J)}{T} &=- \dv{\log(\homosD)}{T}= \dv{\log(\homosD^{-1})}{T}\\
  J &= \frac{1}{\homosD} = \frac{1}{\kappa \sigma_s}
\end{align}
Note that the choice of the initial condition for $J$ is irrelevant, since the multiplicative constant disappears in the ratios.

By summing the particular and homogeneous solutions, we get:
\begin{equation}
\begin{split}
     \sD(T) = - \kappa\sigma_s \sqrt{\frac{45}{4 \pi^3}}\Mpl \int^T_{0} \dd \theta
 \qty[\frac{\sqrt{g_\ast(\theta)}\SigmaD(\theta)}{\qty(1 - \Sigmatstar(\theta))\heff(\theta)\theta^3\tilde H(\theta)\kappa(\theta)\sigma_s(\theta)} ] +\\
 + (C_1 + C_0)\kappa\sigma_s,
\end{split}
\end{equation}
and we find, consistently with the unicity of the solution, that $C_1$ and $C_0$ are not independent integration constants. Finally, by chosing as initial condition $\sD(0)=0$, we find that:
\begin{align}
    \sD(T) = - \sqrt{\frac{45}{4 \pi^3}}\Mpl &\qty(\heff(T) T^3)^\frac{1}{1 - \Sigmatstar(T)} \sigma_s(T) \cdot  \\
    \cdot &\int^T_{0} \dd \theta
 \qty[\frac{\sqrt{g_\ast(\theta)}\SigmaD(\theta)}{\qty(1 - \Sigmatstar(\theta))\tilde H(\theta)\qty(\heff(\theta) \theta^3)^\frac{2- \Sigmatstar(\theta)}{1 - \Sigmatstar(\theta)}\sigma_s(\theta)} ]
\end{align}


\end{document}
